\chapter{Formato dei dati di ingresso}
\label{ch:formato_dati}
\index{Import!formato dati}\index{Excel!schema colonne}\index{CSV!schema colonne}

\grokfig{fig_import_csv_tubature}{File Excel e CSV che entrano
in una rete di tubature pastello, confluiscono in un imbuto e ne
escono righe DataFrame ordinate.}

\epigraph{``Una colonna ben nominata risparmia mezza giornata di
discussioni.''}{Detto di redazione, archivio Borghi, 2025}

\lettrine[lines=3]{Q}{uesto} capitolo specifica i file Excel e CSV
che il sistema accetta. Serve a chi importa per la prima volta
l'anagrafica della propria scuola, a chi prepara un foglio di
prova, a chi riceve un export da un istituto vicino e vuole
sapere se le colonne sono compatibili. La sequenza dei passi
operativi sta nel cap.~\ref{ch:workflow_tipici}, procedura~2; qui
trovi la specifica formale di ogni colonna.

\section{Un endpoint, sette entit\`a}

Tutti gli import passano per uno stesso \emph{endpoint}, cio\`e un
indirizzo HTTP che il browser chiama dietro le quinte quando
premi un bottone della web app:

\begin{lstlisting}
POST /api/import/{entity}
GET  /api/import/{entity}/template
\end{lstlisting}

Il bottone \emph{Importa Excel/CSV} che vedi sopra ogni lista
della web app chiama esattamente quell'endpoint. Le sette
entit\`a accettate sono:

\begin{itemize}
\item \texttt{teachers} -- docenti;
\item \texttt{subjects} -- materie;
\item \texttt{classes} -- classi;
\item \texttt{classrooms} -- aule;
\item \texttt{curricula} -- indirizzi di studio;
\item \texttt{students} -- studenti;
\item \texttt{groups} -- gruppi articolati.
\end{itemize}

Due endpoint accessori, meno frequenti nella pratica, accettano
gli stessi formati: \texttt{assignments} per le cattedre gi\`a
calcolate altrove e \texttt{subject\_group\_weights} per i pesi
materia-gruppo. Il flusso tipico li copre automaticamente dalla
pipeline; servono solo a chi ha gi\`a un foglio Excel preparato.

\section{Convenzioni generali}

Le regole seguenti valgono per tutti gli schemi.

\paragraph{Header.} La prima riga del file contiene i nomi
colonna. Il riconoscimento ignora maiuscole e minuscole; spazi e
underscore si equivalgono (\texttt{Cognome Nome} e
\texttt{cognome\_nome} sono lo stesso campo). Sono accettati alias
italiani e inglesi; ogni schema della sez.~\ref{sec:formato_schemi}
elenca quelli ammessi.

\paragraph{Celle vuote.} Una cella vuota vale \texttt{null}. Una
riga interamente vuota viene ignorata e non compare nei conteggi
del report di import.

\paragraph{Valori booleani.} Sono accettati come veri:
\texttt{true}, \texttt{1}, \texttt{si}, \texttt{vero}, \texttt{x},
\texttt{on} (senza distinzione di maiuscole). Tutto il resto vale
falso, compreso \texttt{0}, \texttt{no} e la cella vuota.

\paragraph{Date.} Formato ISO \texttt{AAAA-MM-GG} (per
esempio \texttt{2010-09-15}). Le date scritte come date native di
Excel funzionano: il sistema legge il valore numerico Excel e lo
converte. Il formato italiano \texttt{15/09/2010} \emph{non} viene
interpretato come data e finisce nella cella come stringa: nei
campi di tipo data si converte in \texttt{null}.

\paragraph{Delimitatore CSV.} Il sistema rileva da solo virgola,
punto e virgola, tabulazione, barra verticale. Il file CSV va
salvato in UTF-8 (con o senza BOM). Gli accenti rovinati in
\texttt{?} o sequenze come \texttt{\`a}\texttt{a} sono il sintomo
classico di un CSV salvato in Windows-1252.

\paragraph{Modalit\`a di import.} Sempre selezionabile dal modal
di import della UI; tre valori:

\begin{description}
\item[\texttt{upsert}] (default) -- aggiorna le righe gi\`a presenti
  (matchate sul campo identificativo) e inserisce le nuove. \`E
  la modalit\`a sicura, la prima da provare sempre.
\item[\texttt{append}] -- come \texttt{upsert} ma ignora gli
  errori delle singole righe e prosegue. Comoda per file di grande
  taglia in cui si tollerano scarti.
\item[\texttt{replace}] -- svuota la tabella prima di importare.
  Cancella anche i record collegati (le associazioni
  classe-materia spariscono insieme alle classi). Usala solo dopo
  aver fatto uno snapshot del database.
\end{description}

\begin{avvertenzabox}
La modalit\`a \texttt{replace} \`e distruttiva e silenziosa: non
chiede conferma una seconda volta. Prima di lanciarla su un
database con vincoli e cattedre gi\`a inseriti, vai sul
Dashboard, card \emph{Import / Export DB}, e premi \emph{Salva
snapshot}.
\end{avvertenzabox}

\section{I sette schemi}
\label{sec:formato_schemi}

Per ciascuna entit\`a riportiamo la tabella dei campi, il campo
identificativo della riga e un mini-esempio inline.

\subsection{Docenti (\texttt{teachers})}

\begin{table}[h]
\centering\small
\begin{tabular}{l l l p{4.5cm}}
\toprule
\textbf{Colonna} & \textbf{Alias italiani} & \textbf{Tipo} & \textbf{Note} \\
\midrule
\texttt{name} & nome, cognome\_nome & str & obbligatorio, identificativo \\
\texttt{matricola} & code & str & matricola interna \\
\texttt{group} & gruppo, classe\_di\_concorso, cdc & str & sigla MIUR (es. A027) \\
\texttt{max\_hours} & max\_ore, ore\_max & int & default 18 \\
\texttt{completion\_hours} & ore\_completamento & int & \\
\texttt{exemption\_hours} & ore\_esonero & int & \\
\texttt{free\_day} & giorno\_libero & str & ``Sabato'', ``Saturday'' \\
\texttt{max\_consecutive} & max\_consecutive\_ore & int & default 5 \\
\texttt{subjects} & materie & csv & ``Mat,Fisica'' \\
\texttt{notes} & note & str & libero \\
\bottomrule
\end{tabular}
\caption{Schema della tabella \texttt{teachers}. La colonna
\texttt{subjects} sostituisce l'insieme materie del docente: se la
ometti, il sistema non tocca le associazioni esistenti; se la
metti, le riscrive.}
\label{tab:formato_teachers}
\end{table}

\paragraph{Esempio.}
\begin{lstlisting}
name,group,max_hours,free_day,subjects
Borghi Giovanni,A027,18,Sabato,"Mat,Fisica"
Bianchi Anna,A012,16,,"Italiano,Latino"
Rossi Luca,A019,12,Lunedi,Storia
\end{lstlisting}

\paragraph{Errori tipici.} Una matricola duplicata su due docenti
con lo stesso \texttt{name} unisce le righe nell'identico record;
una cella \texttt{free\_day} con valore ``sab.'' o ``s.'' resta
inintesa e viene ignorata.

\subsection{Materie (\texttt{subjects})}

\begin{table}[h]
\centering\small
\begin{tabular}{l l l p{5.0cm}}
\toprule
\textbf{Colonna} & \textbf{Alias} & \textbf{Tipo} & \textbf{Note} \\
\midrule
\texttt{name} & nome, materia & str & obbligatorio, identificativo \\
\texttt{pretty\_name} & nome\_esteso & str & per la stampa \\
\texttt{distribute\_days\_weight} & peso\_distribuzione & float & default 0 \\
\texttt{dual\_hours\_weight} & peso\_doppia & float & default 0 \\
\texttt{no\_sixth\_hour\_weight} & peso\_no\_6ora & float & default 0 \\
\texttt{notes} & note & str & libero \\
\bottomrule
\end{tabular}
\caption{Schema di \texttt{subjects}. I tre pesi entrano nella
funzione obiettivo SOFT della Phase B; valori tipici fra 0 e 5.}
\label{tab:formato_subjects}
\end{table}

\paragraph{Esempio.}
\begin{lstlisting}
name,pretty_name,distribute_days_weight,dual_hours_weight
Matematica,Matematica,2.0,1.0
Fisica,Fisica,1.0,2.0
Religione,Religione cattolica,0,0
\end{lstlisting}

\subsection{Classi (\texttt{classes})}

Una classe pu\`o occupare pi\`u righe del file: una riga per coppia
\texttt{(name, subject)}. Il sistema unisce le righe con lo stesso
\texttt{name}.

\begin{table}[h]
\centering\small
\begin{tabular}{l l l p{4.6cm}}
\toprule
\textbf{Colonna} & \textbf{Alias} & \textbf{Tipo} & \textbf{Note} \\
\midrule
\texttt{name} & nome, classe & str & obbligatorio (es. ``1A'') \\
\texttt{year} & anno & int & 1..5 \\
\texttt{section} & sezione & str & ``A'', ``B''\dots \\
\texttt{curriculum} & indirizzo, curriculum\_code & str & codice indirizzo \\
\texttt{n\_students} & studenti & int & numerosit\`a \\
\texttt{subject} & materia & str & una per riga \\
\texttt{hours\_per\_week} & ore, ore\_settimanali & int & ore della materia \\
\texttt{notes} & note & str & libero \\
\bottomrule
\end{tabular}
\caption{Schema di \texttt{classes}. Una classe con cinque materie
occupa cinque righe; i campi anagrafici (\texttt{year},
\texttt{section}, \texttt{curriculum}) bastano sulla prima.}
\label{tab:formato_classes}
\end{table}

\paragraph{Esempio.}
\begin{lstlisting}
name,year,section,curriculum,n_students,subject,hours_per_week
1A,1,A,Scientifico,24,Matematica,5
1A,,,,,Fisica,2
1A,,,,,Italiano,4
1B,1,B,Scientifico,22,Matematica,5
\end{lstlisting}

\paragraph{Identificativo.} \texttt{name}. Se l'indirizzo
\texttt{curriculum} corrisponde al codice di un indirizzo gi\`a in
tabella \texttt{curricula}, il sistema collega anche
\texttt{curriculum\_id}; se non corrisponde, la classe viene creata
ugualmente ma resta scollegata.

\subsection{Aule (\texttt{classrooms})}

\begin{table}[h]
\centering\small
\begin{tabular}{l l l p{4.6cm}}
\toprule
\textbf{Colonna} & \textbf{Alias} & \textbf{Tipo} & \textbf{Note} \\
\midrule
\texttt{name} & nome, aula & str & obbligatorio \\
\texttt{kind} & tipo & str & valori: \texttt{standard}, \texttt{lab\_chimica}, \texttt{lab\_fisica}, \texttt{lab\_informatica}, \texttt{lab\_linguistico}, \texttt{palestra}, \texttt{biblioteca}, \texttt{aula\_speciale} \\
\texttt{capacity} & capienza & int & default 30 \\
\texttt{multi\_class} & multi\_classe & bool & default no \\
\texttt{multi\_class\_max} & multi\_classe\_max & int & default 1 \\
\texttt{notes} & note & str & libero \\
\bottomrule
\end{tabular}
\caption{Schema di \texttt{classrooms}. Il campo \texttt{kind} \`e
ortogonale ai tag; i tag si gestiscono dalla UI dopo l'import.}
\label{tab:formato_classrooms}
\end{table}

\paragraph{Esempio.}
\begin{lstlisting}
name,kind,capacity,multi_class
AulaMagna,aula_speciale,200,true
LabFisica1,lab_fisica,28,false
Palestra,palestra,60,true
A101,standard,28,false
\end{lstlisting}

\subsection{Indirizzi (\texttt{curricula})}

Una riga per coppia \texttt{(code, year, subject)}. Tutte le
righe con lo stesso \texttt{code} si uniscono per formare il
monte ore complessivo dell'indirizzo per ogni anno.

\begin{table}[h]
\centering\small
\begin{tabular}{l l l p{4.6cm}}
\toprule
\textbf{Colonna} & \textbf{Alias} & \textbf{Tipo} & \textbf{Note} \\
\midrule
\texttt{code} & codice, curriculum, indirizzo & str & machine name (``Scientifico'') \\
\texttt{name} & nome & str & display name \\
\texttt{description} & descrizione & str & libero \\
\texttt{score} & punteggio & int & default 1 \\
\texttt{year} & anno & int & 1..5, obbligatorio \\
\texttt{subject} & materia & str & obbligatorio \\
\texttt{hours\_per\_week} & ore, ore\_settimanali & int & ore della materia \\
\bottomrule
\end{tabular}
\caption{Schema di \texttt{curricula}. Una riga per coppia
\texttt{(code, year, subject)}; \texttt{name} e
\texttt{description} bastano sulla prima riga.}
\label{tab:formato_curricula}
\end{table}

\paragraph{Esempio.}
\begin{lstlisting}
code,name,year,subject,hours_per_week
Scientifico,Liceo Scientifico,1,Matematica,5
Scientifico,,1,Fisica,2
Scientifico,,1,Italiano,4
Scientifico,,2,Matematica,5
Scientifico,,2,Fisica,3
\end{lstlisting}

\subsection{Studenti (\texttt{students})}

\begin{table}[h]
\centering\small
\begin{tabular}{l l l p{4.6cm}}
\toprule
\textbf{Colonna} & \textbf{Alias} & \textbf{Tipo} & \textbf{Note} \\
\midrule
\texttt{last\_name} & cognome & str & obbligatorio \\
\texttt{first\_name} & nome & str & obbligatorio \\
\texttt{birth\_date} & data\_nascita & date & ISO \texttt{AAAA-MM-GG} \\
\texttt{gender} & sesso & str & M, F, other \\
\texttt{email} & & str & \\
\texttt{student\_code} & matricola, codice & str & id esterno \\
\texttt{class\_name} & classe, class & str & ``1A'' \\
\texttt{notes} & note & str & libero \\
\bottomrule
\end{tabular}
\caption{Schema di \texttt{students}. L'identificativo \`e la
tripletta \texttt{(last\_name, first\_name, birth\_date)};
\texttt{student\_code} \`e opzionale ma utile per riferimenti
incrociati con altri sistemi.}
\label{tab:formato_students}
\end{table}

\paragraph{Esempio.}
\begin{lstlisting}
last_name,first_name,birth_date,class_name,student_code
Rossi,Mario,2010-09-15,1A,STU0001
Bianchi,Anna,2010-11-03,1A,STU0002
Verdi,Luca,2010-04-22,1B,STU0003
\end{lstlisting}

\subsection{Gruppi articolati (\texttt{groups})}

Una riga per coppia \texttt{(name, subject)}. Gli studenti si
elencano nella colonna \texttt{student\_codes} come lista
virgola-separata.

\begin{table}[h]
\centering\small
\begin{tabular}{l l l p{4.6cm}}
\toprule
\textbf{Colonna} & \textbf{Alias} & \textbf{Tipo} & \textbf{Note} \\
\midrule
\texttt{name} & nome, group, gruppo & str & obbligatorio \\
\texttt{kind} & tipo & str & \texttt{splitting}, \texttt{language}, \texttt{religion}, \texttt{support} \\
\texttt{description} & descrizione & str & libero \\
\texttt{subject} & materia & str & materia del gruppo \\
\texttt{hours\_per\_week} & ore & int & ore alla settimana \\
\texttt{student\_codes} & matricole, students & csv & codici virgola-separati \\
\texttt{notes} & note & str & libero \\
\bottomrule
\end{tabular}
\caption{Schema di \texttt{groups}. Gli studenti referenziati per
\texttt{student\_code} devono essere gi\`a presenti in
\texttt{students}; quelli non trovati finiscono fra gli
\texttt{errors} del report ma non bloccano il resto dell'import.}
\label{tab:formato_groups}
\end{table}

\paragraph{Esempio.}
\begin{lstlisting}
name,kind,subject,hours_per_week,student_codes
AltRel_1A,religion,Religione,1,"STU0001,STU0002,STU0003"
Ing_Avanzato,language,Inglese,3,"STU0010,STU0011,STU0014"
\end{lstlisting}

\section{Schema dei file di vincoli}

I vincoli logici hanno il loro file di import, separato dagli
schemi anagrafici. L'endpoint \`e
\texttt{POST /api/dashboard/constraints/import-file}; accetta JSON
o xlsx con record nella forma seguente.

\begin{table}[h]
\centering\small
\begin{tabular}{l l p{6.2cm}}
\toprule
\textbf{Campo} & \textbf{Tipo} & \textbf{Note} \\
\midrule
\texttt{kind} & str & obbligatorio. ``unavailability\_matrix'', ``logical'', ``coteaching'', ``room\_preference''\dots \\
\texttt{scope} & str & obbligatorio. ``teacher'', ``class'', ``classroom'', ``global'' \\
\texttt{owner\_name} & str & nome del docente/classe/aula a cui appartiene \\
\texttt{level} & str & ``hard'', ``soft'', ``preferred''. Default ``hard'' \\
\texttt{expression} & str & espressione DSL del vincolo \\
\texttt{soft\_penalty} & int & peso se \texttt{level=soft} \\
\texttt{description} & str & etichetta per il pannello UI \\
\bottomrule
\end{tabular}
\caption{Schema dei record di vincolo. Il sistema risolve
\texttt{owner\_name} contro l'anagrafica corrispondente a
\texttt{scope}; un nome non trovato finisce negli \texttt{errors}
del report senza bloccare il resto.}
\label{tab:formato_vincoli}
\end{table}

\paragraph{Esempio JSON.}
\begin{lstlisting}
[
  {
    "kind": "logical",
    "scope": "teacher",
    "owner_name": "Bianchi Anna",
    "level": "hard",
    "expression": "mai sab",
    "description": "Bianchi non lavora il sabato"
  },
  {
    "kind": "logical",
    "scope": "class",
    "owner_name": "5A",
    "level": "soft",
    "soft_penalty": 5,
    "expression": "evita ven6",
    "description": "5A evita l'ultima ora del venerdi"
  }
]
\end{lstlisting}

\paragraph{Scorciatoia DSL.} Un'espressione con un giorno senza ora
(\texttt{mai sab}, \texttt{lun OR mar}) viene espansa automaticamente
nella forma esplicita che il parser richiede, cio\`e
\texttt{(sab8 OR sab9 OR sab10 OR sab11 OR sab12 OR sab13)}. Lo
puoi usare anche nei file scritti a mano.

\section{La risposta dell'endpoint: \texttt{ImportReport}}

Ogni import restituisce un oggetto JSON con la stessa struttura:

\begin{lstlisting}
{
  "ok": true,
  "entity": "teachers",
  "n_inserted": 12,
  "n_updated": 3,
  "n_skipped": 1,
  "n_total_rows": 16,
  "messages": [],
  "errors": ["riga 14: name mancante"]
}
\end{lstlisting}

Nella web app il bottone \emph{Importa} apre un piccolo modal che
riassume i numeri e, in caso di errori, ne mostra l'elenco
scorrendo. Il numero significativo per capire se l'import \`e
andato bene \`e \texttt{n\_total\_rows} confrontato con la somma
\texttt{n\_inserted + n\_updated}. Se i due numeri coincidono,
nessuna riga \`e stata scartata.

\begin{table}[h]
\centering\small
\begin{tabular}{l l p{6.0cm}}
\toprule
\textbf{Campo} & \textbf{Tipo} & \textbf{Significato} \\
\midrule
\texttt{ok} & bool & l'import nel suo complesso \`e riuscito \\
\texttt{entity} & str & il nome dell'entit\`a importata \\
\texttt{n\_inserted} & int & righe nuove create \\
\texttt{n\_updated} & int & righe esistenti aggiornate \\
\texttt{n\_skipped} & int & righe ignorate o scartate \\
\texttt{n\_total\_rows} & int & righe del file (esclusa la header) \\
\texttt{messages} & list & note informative non bloccanti \\
\texttt{errors} & list & dettaglio degli errori per riga \\
\bottomrule
\end{tabular}
\caption{Campi del report di import.}
\label{tab:formato_report}
\end{table}

\section{Errori tipici e come correggerli}

\begin{erroricomunibox}
\textbf{``Tutti i record vanno in \texttt{n\_skipped}''}. Le
intestazioni di colonna non sono state riconosciute. Verifica che
i nomi colonna corrispondano a quelli dello schema o ai loro
alias. Apri il template scaricato (\emph{Scarica template} sopra
la lista) e copia le intestazioni una a una.

\textbf{``Le date sono diventate \texttt{null}''}. Probabilmente
nel foglio Excel erano scritte in formato italiano
\texttt{15/09/2010}. Convertile in date Excel (formato cella
\emph{Data}) o riscrivile in formato ISO
\texttt{AAAA-MM-GG} \texttt{2010-09-15}.

\textbf{``Gli accenti sono rovinati''}. Salva il CSV in formato
\emph{CSV UTF-8 (delimited)}. In Excel italiano questa opzione
si trova in \emph{File} $\to$ \emph{Salva con nome} $\to$
selettore tipo file. Il formato \emph{CSV (separato dalla
virgola)} produce file Windows-1252 che il sistema non legge
correttamente sugli accenti.

\textbf{``Importo \texttt{students} con \texttt{class\_name}
ma la classe non viene trovata''}. La classe \texttt{1A} deve
esistere in \texttt{classes} \emph{prima} dell'import degli
studenti. L'ordine raccomandato \`e:
\texttt{curricula} $\to$ \texttt{subjects} $\to$ \texttt{teachers}
$\to$ \texttt{classes} $\to$ \texttt{classrooms} $\to$
\texttt{students} $\to$ \texttt{groups}.

\textbf{``Il file ha 50 righe ma \texttt{n\_total\_rows} ne riporta
solo 30''}. Probabilmente venti righe avevano tutte le celle
vuote (rimaste dopo un Trova/Sostituisci che ha cancellato un
campo obbligatorio). Le righe interamente vuote vengono ignorate
silenziosamente.

\textbf{``L'import in modalit\`a \texttt{replace} mi ha cancellato
le cattedre''}. Comportamento previsto: \texttt{replace} svuota
anche i record collegati (le associazioni classe-materia spariscono
insieme alle classi). Usa \texttt{upsert} se vuoi solo aggiornare
i campi delle classi senza toccare le cattedre.

\textbf{``Il bottone \emph{Importa} resta grigio''}. Probabilmente
il file pesa pi\`u del limite (\SI{16}{MB} di default). Spezzalo in
due o aumenta il limite nella variabile d'ambiente
\texttt{TIMETABLE\_MAX\_UPLOAD\_MB} prima di riavviare il
backend.
\end{erroricomunibox}

\fleuron

\begin{biblioparvabox}
Lo schema canonico di import \`e mantenuto in
\texttt{webui/docs/import\_format.md} dentro il repository del
progetto. Quel file e questo capitolo si tengono in pari per
costruzione: chi modifica uno ha l'obbligo morale di aggiornare
anche l'altro. Le coppie alias italiani/inglesi sono raccolte
nel modulo \texttt{webui/backend/routers/imports.py},
funzione \texttt{\_alias\_map} per ogni entit\`a.
\end{biblioparvabox}
